% Borrador del capitulo del marco geologico
\chapter{Marco geológico}\label{cap:mg}
El Complejo Volcánico Los Humeros (CVLH) se localiza en los límites entre los estados de Puebla y Veracruz. Esta región constituye una de las áreas volcánicas más activas del país y ha sido objeto de numerosos estudios debido a su compleja evolución eruptiva y a su relevancia como campo geotérmico de alta importancia estratégica en los recursos enérgeticos de México.

\section{Contexto regional}
El CVLH está ubicado en el sector oriental del FVTM, un arco continental de ~1000 km de largo activo desde el Mioceno, formado por la subducción de las placas de Cocos y Rivera bajo la placa de Norteamérica desde el Neógeno \autocite{Bloque1}. Este sector oriental se superpone a una corteza continental gruesa (45-55 km) de basamento Precámbrico a Paleozoico y se caracteriza por una compleja historia volcánica \autocite{Bloque1}.

Específicamente, el CVLH es el volcán más septentrional de la Cuenca de Serdán-Oriental, una cuenca cerrada con vulcanismo bimodal (basáltico y riolítico) durante el Plioceno-Holoceno. El basamento sobre el cual yace el CVLH es complejo:

\begin{itemize}
    \item \textbf{Basamento cristalino:} En la base se encuentra el Macizo de Teziutlán, un complejo de rocas metamórficas e intrusivas del Paleozoico.
    \item \textbf{Secuencia sedimentaria mesozoica:} Sobre el basamento cristalino se encuentra una potente secuencia ($\leq$ 3000m) de calizas y lutitas, intensamente plegada y fallada durante la Orogenia Laramide (Cretácio Tardío - Eoceno).
    \item \textbf{Intrusiones Terciarias:} Cuerpos graníticos y sieníticos del Eoceno-Oligoceno intrusionaron la secuencia sedimentaria.
    \item \textbf{Vulcanismo pre-CVLH:} La actividad volcánica regional inició en el Mioceno - Plioceno con la efusión de lavas andesíticas y basaltos-andesíticos (Formaciones Teziutlán, Cuyoaco y Alseseca) fechadas entre 1.46 y 10.5 Ma. Estas rocas andesíticas hoy constituyen la principal roca reservorio del campo geotérmico.
\end{itemize}

\section{Evolución volcánica y estratigráfica}
La evolución del CVLH ha sido reevaluada significativamente gracias a nuevas dataciones de alta precisión (U/Th y 40Ar/39Ar), que demuestran que es un sistema mucho más joven de lo que se pensaba. La actividad se puede dividir en tres etapas principales:
\begin{itemize}
    \item \textbf{Estapa pre-caldera (>164 ka):} Inició con la extrusión de domos y flujos de lava de composición riolítica en los flancos oeste y sur de la futura caldera. Las nuevas dataciones ubican esta actividad entre 693 ± 1.9 ka y 270 ± 17 ka. Estos domos, anteriormente considerados post-caldera, ahora se clasifican como pre-caldera gracias a la nueva y más joven edad de la primera gran erupción.
    \item \textbf{Etapa de caldera (164 - 69 ka):} Caracterizada por eventos explosivos masivos.
    \begin{itemize}
        \item \textbf{Colapso de la caldera Los Humeros (~164 ka):}  Fue la erupción individual más voluminosa registrada en el CVTM, con un volumen reevaluado de ~290 km³ de Roca Densa Equivalente (DRE). Este evento formó la caldera principal Los Humeros (21 x 15 km) y depositó la Ignimbrita Xáltipan, una unidad clave que actúa como roca sello del reservorio geotérmico. Nuevas dataciones precisas establecen la edad de este evento en 164 ± 4.2 ka, mucho más joven que los 460 ka previamente estimados.
        \item \textbf{Erupciones plinianas - Toba Faby (~70 ka):} Una serie de erupciones explosivas depositaron la Toba Faby, una secuencia de capas de caída de pómez riodacítica a andesítica. La edad de este evento ha sido redefinida a 70 ± 23 ka.
        \item \textbf{Colapso de la caldera Los Potreros (~69 ka):}  Un segundo evento formador de caldera, de menor volumen (~15 km³ DRE), depositó la Ignimbrita Zaragoza y formó la caldera anidada Los Potreros (9-10 km de diámetro). Su edad ha sido revisada a 69 ± 16 ka.
    \end{itemize}

    \item \textbf{Etapa post-caldera (<69 ka al Holoceno):} 
    \begin{itemize}
        \item \textbf{Fase resurgente (Pleistoceno Tardío):} Incluyó el emplazamiento de domos riolíticos y dacíticos en el interior de la caldera (~45 ka) y erupciones explosivas que formaron las Tobas Xoxoctic (~50 ka) y Llano (~28 ka).
        \item \textbf{Fase de vulcanismo bimodal y de fractura anular (Holoceno):} Es la actividad más reciente, con erupciones de composición bimodal y flujos de lava basálticos-andesíticos a traquíticos a lo largo de fracturas anulares. Destacan la Toba Cuicuiltic (~7.3 ka), producto de una erupción bimodal simultánea, y un extenso campo de lavas en el borde sur emplazado entre ~7.3 ka y 2.86 ka. La actividad más reciente son flujos de basalto de olivino (~4 ka).
    \end{itemize}
    
\end{itemize}
\section{Litología y petrografía del CVLH y su subsuelo}
El CVLH presenta una amplia gama de litologías, desde basaltos hasta riolitas de alta sílice. La estratigrafía del subsuelo, revelada por más de 50 pozos geotérmicos, muestra una gran heterogeneidad. Las principales unidades son:
\begin{itemize}
    \item \textbf{Basamento pre-volcánico:} Calizas y lutitas mesozoicas metamorfizadas a hornfels y skarn, con intrusiones de dioritas y granodioritas.
    \item \textbf{Grupo Pre-caldera:}
    \begin{itemize}
        \item \textbf{Andesitas basales (Mioceno):} Andesitas con hornblende y basaltos en la base de algunos pozos, correlacionadas con el vulcanismo de Cerro Grande (8.5-11 Ma)
        \item \textbf{Andesitas con piroxeno (Plioceno - Pleistoceno):} Es la sucesión más potente (>1500 m) y constituye la principal roca reservorio del sistema geotérmico. Su edad se ha determinado entre 1.46 y 2.61 Ma.
        \item \textbf{Riolitas y dacitas superiores:} Domos y flujos de lava intercalados con las andesitas, con distribución local.
    \end{itemize}

    \item \textbf{Grupo Caldera:}
    \begin{itemize}
        \item \textbf{Ignimbrita Xáltipan:} Es la unidad que actúa como roca sello del reservorio. Presenta variaciones significativas de soldamiento, desde no soldada a altamente soldada (lito-facies tipo lava-like), controlando fuertemente la permeabilidad. Alcanza espesores de hasta 880 m dentro de la caldera.
        \item \textbf{Toba Faby e Ignimbrita Zaragoza:} Depósitos piroclásticos intercalados con lavas, presentes en el subsuelo pero con distribución y espesores variables.
    \end{itemize}

    \item \textbf{Grupo Post-caldera:} Lavas basálticas y andesíticas, y depósitos de escoria y pómez en los niveles más someros de los pozos.
\end{itemize}

\section{Estructuras geológicas y tectónica activa}
La geología estructural del CVLH es el resultado de la interacción entre estructuras tectónicas heredadas y deformaciones volcano-tectónicas recientes.

\begin{itemize}
    \item \textbf{Estructuras heredadas:} El basamento está afectado por pliegues y cabalgaduras de la Orogenia Laramide con orientación NW-SE, y por fallas normales NE-SW del Eoceno-Plioceno. Estas estructuras preexistentes han controlado la ubicación del sistema magmático, los colapsos de caldera y las deformaciones posteriores.
    \item \textbf{Estructuras de colapso:} Se han formado dos calderas anidadas. La Caldera Los Humeros (21x15 km) es una estructura asimétrica tipo \textit{trap-door} con el mayor colapso en el sector SO. La Caldera Los Potreros (9x10 km) está anidada en su interior. Los bordes de ambas están definidos por escarpes morfológicos y fallas anulares con echado hacia afuera de la caldera.
    \item \textbf{Resurgencia y fallamiento activo:} El piso de la caldera Los Potreros está afectado por un sistema de fallas activas relacionadas con un proceso de resurgencia. Este sistema está dominado por un enjambre de fallas principal con orientación NNW-SSE (Fallas Maxtaloya-Los Humeros) y fallas secundarias N-S, NE-SW y E-W. Datos de sismicidad y análisis cinemáticos indican un régimen de esfuerzos complejo, con fallamiento normal en la parte superior del reservorio (<1.5 km) y fallamiento inverso en la parte inferior, sugiriendo una fuente de presión local (magmática/hidrotermal) que causa el levantamiento del piso de la caldera.
\end{itemize}

\section{Geología reciente y procesos activos}
El CVLH es un sistema volcánico activo, con erupciones registradas hasta el Holoceno tardío.

\begin{itemize}
    \item \textbf{Vulcanismo Holoceno:} La actividad más reciente (últimos 10,000 años) se concentra en el borde sur y en el interior de la caldera. Destaca la erupción del Miembro Cuicuiltic (~7.3 ka), un evento bimodal que involucró la activación simultánea de al menos tres centros eruptivos distintos, uno central de composición traquidacítica (Pliniano) y otros fisurales de composición basáltica-andesítica (Estromboliano).
    \item \textbf{Campo de lavas del borde sur:}  La actividad efusiva más reciente está representada por seis flujos de lava principales (Victoria, Texcal, Tepeyahualco, Sarabia, El Limón y El Pájaro) emplazados entre ~7.3 ka y 2.86 ka. Este campo de lavas muestra una gran diversidad composicional (basalto a traquita) y fue generado en tres etapas eruptivas principales.
    \item \textbf{Sistema magmático actual:} La diversidad de productos y la ocurrencia de erupciones bimodales simultáneas sugieren la existencia de un sistema magmático complejo y heterogéneo, posiblemente compuesto por múltiples reservorios interconectados o lentes de magma dentro de una masa de cristales parcialmente consolidada (crystal mush). La presencia de vulcanismo basáltico muy joven (~4 ka) indica recargas recientes de magma en el sistema.
\end{itemize}

\section{Implicaciones geotérmicas}
La reevaluación geológica del CVLH tiene implicaciones fundamentales para la exploración y explotación de sus recursos geotérmicos. Las nuevas edades, en particular la de la erupción Xáltipan (164 ka vs. 460 ka), indican que el sistema magmático es mucho más joven de lo que se pensaba. Esto, sumado al gran volumen magmático estimado (>1150 km³) y a la evidencia de recargas magmáticas holocénicas, aumenta significativamente el potencial geotérmico al implicar una fuente de calor más robusta y duradera.

El reservorio geotérmico principal se localiza en la secuencia de andesitas pre-caldera (Fm. Teziutlán) a profundidades promedio de 1400 m. Este reservorio está sellado por la Ignimbrita Xáltipan, cuya permeabilidad es muy variable y controlada por su grado de soldamiento. Se han identificado al menos dos niveles de reservorio: uno somero, dominado por líquido (~1025-1600 m s.n.m.), y otro más profundo, rico en vapor y con baja saturación de líquido (~100-850 m s.n.m.).

La permeabilidad en el reservorio es principalmente secundaria y está controlada por el sistema de fallas de resurgencia activo (fallas Maxtaloya-Los Humeros) y fracturas asociadas. El flujo de fluidos hidrotermales asciende a través de estas fallas subverticales. La alteración hidrotermal propilítica, con minerales como epidota y magnetita secundaria, está asociada a estas zonas de flujo, creando contrastes geofísicos (resistividad y susceptibilidad magnética) que son clave para la exploración.

Junto con las altas temperaturas registradas (>380°C), la juventud del sistema magmático y la evidencia de una fuente de calor a profundidades de 5-7 km hacen del CVLH un objetivo prometedor para el desarrollo de sistemas geotérmicos supercalientes.
